\subsection{轨道}

命题 14. 设 \(\mathbf{G}\) 为李群，\(\mathbf{X}\) 为解析流形，且 \((g,x)\mapsto gx\) 是 \(\mathbf{G}\) 在 \(\mathbf{X}\) 上的解析左作用律。设 \(x\in \mathbf{X}\)。假设相应的轨道映射 \(\rho (x)\) 是次浸入（如果 \(\mathbf{K}\) 的特征为 0 且 \(\mathbf{X}\) 是有限维的，这总是如此（命题 9 的推论））。令 \(\mathbf{G}_{\mathbf{x}}\) 为稳定子 \(x\) 在 \(\mathbf{G}\) 中的子群。

\begin{enumerate}
\def\labelenumi{(\roman{enumi})}
\item
  \(\mathrm{G}_x\) 是李子群，且 \(\mathrm{T}_e(\mathrm{G}_x) = \mathrm{Ker}\, \mathrm{T}_e(\rho(x))\)。
\item
  齐性空间的典范映射 \(i_x\) 从 \(G / G_x\) 到 \(X\) 是浸入，其像为 \(Gx\)。
\item
  如果进一步轨道 \(\mathbf{G}x\) 是局部闭的且 \(\mathbf{G}\) 的拓扑具有可数基，则 \(\mathbf{G}x\) 是 \(\mathbf{X}\) 的子流形，\(i_x\) 是从流形 \(\mathrm{G} / \mathrm{G}_x\) 到流形 \(\mathbf{G}x\) 的同构，且 \(\mathrm{T}_x(\mathbf{G}x) = \operatorname{Im} \mathrm{T}_e(\rho(x))\)。
\end{enumerate}

\(x\) 在 \(\rho(x)\) 下的逆像是 \(G_x\)。由于 \(\rho(x)\) 是次浸入，\(G_x\) 是子流形，并且对于所有 \(g \in G\)，其切空间 \(J\) 是 \(gG_x = \rho(x)^{-1}(gx)\) 在 \(g\) 处的切空间，且等于 \(\operatorname{Ker} T_g(\rho(x))\)（Differentiable and Analytic Manifolds, R, 5.10.5），从而得到 (i)。令 \(\pi: G \to G/G_x\) 为典范映射。于是 \(i_x \circ \pi = \rho(x)\)。由于 \(G/G_x\) 是 \(G\) 的商流形，这个等式表明 \(i_x\) 是解析的。此外，\(T_g(\rho(x))\) 和 \(T_g(\pi)\) 的核都等于 \(J\)。因此 \(T_{\pi(g)}(i_x)\) 是单射。\(T_{\pi(g)}(i_x)\) 的像等于 \(T_g(\rho(x))\) 的像，因而有拓扑补。这证明了 (ii)。

假设 \(Gx\) 是局部闭的。于是 \(Gx\) 中每一点在 \(Gx\) 中都有一个邻域，它同胚于完备度量空间的闭子空间，因而是 Baire 空间。因此 \(Gx\) 是 Baire 空间（General Topology, Chapter IX, § 5, Proposition 4）。如果 \(G\) 具有可数基，则 \(i_x\) 是从 \(G / G_x\) 到 \(Gx\) 的同胚（General Topology, Chapter IX, § 5）。于是由 (ii) 和 Differentiable and Analytic Manifolds, R, 5.8.3，\(i_x\) 是从流形 \(G / G_x\) 到流形 \(Gx\) 的同构，并且

\[
\mathrm{T} _ {x} (\mathrm{G} x) = \operatorname{Im} \mathrm{T} _ {\pi (e)} (i _ {x}) = \operatorname{Im} \mathrm{T} _ {e} (\rho (x)).
\]

备注。设 G 为有限维李群，X 为 \(C^{r}\) 类流形，且 \((g, x) \mapsto gx\) 是 G 在 X 上的 \(C^{r}\) 类左作用律。则命题 14 仍然成立。唯一需要不同证明的点是 \(G_{x}\) 为李子群这一事实。但若 \(r \neq \omega\)，则 K = R；由于 \(G_{x}\) 显然是闭的，\(G_{x}\) 由§8、定理 2 得出是李子群。

推论。设 G 为拓扑具有可数基的李群，X 为

非空有限维解析流形，且 G 在 X 上有解析左作用律。假设 G 在 X 上传递作用且 K 的特征为 0。则 X 是 G 的李齐性空间。

令 \(x \in X\)。\(x\) 的轨道等于 \(X\)，是闭的，因此可以应用命题 14 (iii)。

